\documentclass[10pt,a4paper]{amsart}
\usepackage[margin=19mm]{geometry}
\usepackage{amsmath,amssymb,mathtools}
\usepackage[scr=rsfs,cal=euler]{mathalfa}
\usepackage{xfrac}
\usepackage[unicode,hidelinks,bookmarks=false]{hyperref}
\newcommand{\ie}{i.e.\ }
\newcommand{\cf}{cf.\ }
\newcommand{\resp}{respectively\ }
\newcommand{\Subsection}{\subsection}
\providecommand{\nobreakdash}{\nobreak\hbox{-}}
\setlength{\emergencystretch}{2em}
\multlinegap0pt
\usepackage{bm}
\usepackage{bbm}
\usepackage{dsfont}
\usepackage{xspace}
\usepackage{cancel}
\usepackage{mathtools}
\usepackage{cleveref}
\usepackage{amsthm}
\makeatletter
\@ifundefined{theorem}{\newtheorem{theorem}{Theorem}}{}
\@ifundefined{corollary}{\newtheorem{corollary}[theorem]{Corollary}}{}
\@ifundefined{lemma}{\newtheorem{lemma}[theorem]{Lemma}}{}
\makeatother
\DeclareMathOperator{\Star}{Star}
\DeclareMathOperator{\Trop}{Trop}
\DeclareMathOperator{\rank}{rank}
\newcommand{\EE}{\mathbb{E}}
\newcommand{\GG}{\mathbb{G}}
\newcommand{\HH}{\mathbb{H}}
\newcommand{\RR}{\mathbb{R}}
\newcommand{\bbone}{\mathbbm{1}}
\newcommand{\ee}{\mathbbm{e}}
\title[The tropical galaxy of a Laman graph]{The tropical galaxy of a Laman graph}
\author{Amelia Bielby, Arushi Chauhan, Cassia Pearce, Yue Ren}
\date{Source: arXiv:2511.09246v1 (CC BY 4.0)}
\begin{document}
\maketitle

\noindent\emph{Source and licence.} The tropical galaxy of a Laman graph. Version arXiv:2511.09246v1 (CC BY 4.0), distributed under the Creative Commons Attribution 4.0 International licence, as recorded in the arXiv licence metadata for this exact version and shown on the abstract page https://arxiv.org/abs/2511.09246v1. Full rights notice: licenses/arxiv-2511.09246-CC-BY-4.0.txt.\par\smallskip

\noindent\textit{Editorial scope.} Counted unit: the Lemma whose source label is \texttt{lem:excisionToStar} in the article (source0.tex lines 1121--1132, source label \texttt{lem:excisionToStar}), delivered together with the article's own complete original proof of it (source0.tex lines 1133--1188, one \texttt{proof} environment of 56 lines). No mathematics is written, repaired or re-derived: statement and proof are the article's own text line for line. Required context retained uncounted: lem:excisionsCommuting (lines 1083--1089); cor:excisionsAndChains (lines 1110--1113); cor:excisionsAndChains2 (lines 1115--1118). Every definition, display and label of those lines is retained unchanged. Omitted from this package and not redistributed: 1--1082, 1090--1109, 1114--1114, 1119--1120, 1189--2094. Nothing inside a retained excerpt is omitted or altered: every retained body is delivered exactly as the article's lines are written, so no omission receipt of any kind accompanies this package. Citations: the retained excerpts carry no citation command at all, so no bibliography entry of the article is transcribed and no reference is redistributed. Reference adaptation: none was needed and none was made; every reference inside the delivered unit resolves inside it, so no unresolved reference or citation is produced. Printed slips kept literal: no printed word, symbol, hypothesis or number of the counted statement or of its proof was altered. Layout and class: the article's own class is \texttt{amsart} with packages including inputenc, biblatex, hyperref, amsmath, amsthm, thmtools, amssymb, amsfonts, bbm, mathtools, stmaryrd, enumitem, mathrsfs, comment; the wrapper is the lane's pinned \texttt{amsart} editorial wrapper at 10pt on A4, which restates the theorem-like environments the retained text needs and adds the article's own macro definitions that the retained text needs (\texttt{\textbackslash EE}, \texttt{\textbackslash GG}, \texttt{\textbackslash HH}, \texttt{\textbackslash RR}, \texttt{\textbackslash Star}, \texttt{\textbackslash Trop}, \texttt{\textbackslash bbone}, \texttt{\textbackslash ee}, \texttt{\textbackslash rank}), with extra wrapper packages (bm, bbm, dsfont, xspace, cancel, mathtools, cleveref). The delivered numbering is the unit's own. Assets: no retained span contains a figure environment, a graphics-inclusion command, a PDF-page inclusion, a table or a TikZ picture; no image file is copied into this package, the package declares no asset dependency and no image file is redistributed with it. Licence: version arXiv:2511.09246v1 of this article is distributed under the Creative Commons Attribution 4.0 International licence, as recorded in the arXiv licence metadata for this exact version and shown on the abstract page https://arxiv.org/abs/2511.09246v1. A full rights notice is delivered at licenses/arxiv-2511.09246-CC-BY-4.0.txt. Characters: every retained excerpt is pure ASCII, so the delivered engine, XeLaTeX with its default Latin Modern text font, reports no missing character.
\begin{lemma}
  \label{lem:excisionsCommuting}
  Let $\GG$ be a multigraph, and let $\ee_1,\ee_2\in\EE(\GG)$ be two of its multiedges on two different connected components.  Then
  \begin{equation*}
    \GG\curvearrowright\ee_1\curvearrowright\ee_2 = \GG\curvearrowright\ee_2\curvearrowright\ee_1.
  \end{equation*}
\end{lemma}

\begin{corollary}
  \label{cor:excisionsAndChains}
  Let $\GG$ be a multigraph.  Then any chain of excisions $\GG\curvearrowright\ee_1\curvearrowright\dots\curvearrowright\ee_r$ gives rise to a chain of flats $\emptyset = F_0\subsetneq \ee_1 \subsetneq \ee_1\cup \ee_2 \subsetneq\dots\subsetneq \ee_1\cup\dots\cup\ee_r$ with $\rank(\ee_1\cup\dots\cup\ee_j)=j$, and any chain of flats $\emptyset=F_0\subsetneq F_1\subsetneq\dots\subsetneq F_r$ with $\rank(F_j)=j$ gives rise to a chain of excisions $\GG\curvearrowright F_1\curvearrowright F_2\setminus F_1\curvearrowright \dots\curvearrowright F_r\setminus F_{r-1}$.
\end{corollary}

\begin{corollary}
  \label{cor:excisionsAndChains2}
  Let $\GG$ be a multigraph, and fix a chain of excisions $\HH\coloneqq \GG\curvearrowright\ee_1\curvearrowright\dots\curvearrowright\ee_r$ and its corresponding chain of flats $F_\bullet\colon \emptyset=F_0\subsetneq F_1\subsetneq\dots\subsetneq F_r$ from \cref{cor:excisionsAndChains}.  Then any chain of flats $F_\bullet'$ on $\HH$ beginning with $F_{i}'=\ee_i$ is a chain of flats on $\GG$ extending $F_\bullet$, and vice versa.
\end{corollary}

\begin{lemma}
  \label{lem:excisionToStar}
  Let $\GG$ be a multigraph and let $\HH\coloneqq \GG\curvearrowright\ee_1\curvearrowright\dots\curvearrowright\ee_r$ arise from a chain of excisions.  Let $\GG_1,\dots,\GG_s$ be the connected components of $\GG$.  By \cref{lem:excisionsCommuting}, we may reorder $\ee_1,\dots,\ee_r$ so that $\EE(\GG_k)\supseteq \{\ee_{j_{k-1}+1},\dots,\ee_{j_k}\}$ for $0=j_0<\dots<j_s=r$. Then the following is a valid cone in $\Trop(\GG)$:
  \begin{align*}
    \sigma \coloneqq &\RR_{\geq 0}\cdot \bbone_{\ee_1}+\RR_{\geq 0}\cdot \bbone_{\ee_1\cup\ee_2}+\dots+\RR_{\geq 0}\cdot \bbone_{\ee_1\cup\dots\cup\ee_{j_1}}\\
    &\quad +\RR_{\geq 0}\cdot \bbone_{\ee_{j_1+1}}+\RR_{\geq 0}\cdot \bbone_{\ee_{j_1+1}\cup\ee_{j_1+2}}+\dots+\RR_{\geq 0}\cdot \bbone_{\ee_{j_1+1}\cup\dots\cup\ee_{j_2}}\\
    &\hspace{4.75mm} \vdots\\
    &\quad +\RR_{\geq 0}\cdot \bbone_{\ee_{j_{s-1}+1}}+\RR_{\geq 0}\cdot \bbone_{\ee_{j_{s-1}+1}\cup\ee_{j_{s-1}+2}}+\dots+\RR_{\geq 0}\cdot \bbone_{\ee_{j_{s-1}+1}\cup\dots\cup\ee_{j_s}}\\
    &\quad +\RR\cdot\bbone_{[m]}\in \Trop(\GG).
  \end{align*}
  Moreover, $\Trop(\HH)=\Star_{\Trop(\GG)}(\sigma)$.
\end{lemma}

\begin{proof}
  We may assume without loss of generality that $\GG$ is connected, which means that the only connected components of $\HH$ are $\HH_1,\dots,\HH_r$ with $E(\HH_j)=\ee_j$, and a component $\HH'\subseteq\HH$ with $E(\HH')=E(\GG)\setminus E$ for $E\coloneqq \ee_1\cup\dots\cup\ee_r$.  By \cref{cor:excisionsAndChains}, the following is a valid chain of flats on $\GG$:
  \begin{equation*}
    F_{E,\bullet}\colon\;\; \emptyset\subsetneq \ee_1\subsetneq \ee_1\cup\ee_2\subsetneq\dots\subsetneq \ee_1\cup\dots\cup\ee_r \subseteq E.
  \end{equation*}
  Hence $\sigma=\sigma(F_{E,\bullet})\in\Trop(\GG)$.  It remains to show that $\Trop(\HH)=\Star_{\Trop(\GG)}(\sigma)$.
  
  For the ``$\subseteq$'' inclusion, let $\tau\in\Trop(\HH)$ be a maximal cone.  By definition, $\tau$ arises from a maximal chain of flats on each connected component of $\HH$, which we can rearrange to become a single maximal chain of flats on $\HH$ beginning with the $\ee_i$:
  \begin{equation*}
    F_{[m],\bullet}\colon\;\; \emptyset\subsetneq \ee_1\subsetneq \dots\subsetneq \ee_1\cup\dots\cup\ee_r\subsetneq E\cup F_1\subsetneq E\cup F_2\subsetneq\dots\subsetneq E\cup F_\ell\subsetneq [m].
  \end{equation*}

  By \cref{cor:excisionsAndChains2}, this is also a valid maximal chain of flats on $\GG$ which corresponds to a maximal cone $\sigma(F_{[m],\bullet})\in\Trop(\GG)$.  We now claim that $\tau=\Star_{\sigma(F_{[m],\bullet})}(v)$ for any $v\in\mathrm{Relint}(\sigma)$, say $v = \bbone_{\ee_1}+\dots+\bbone_{\ee_1\cup\dots\cup\ee_r}$.

  Let $u=(u_i)_{i\in[m]}\in \mathrm{Relint}(\tau)$, i.e., $(u_i)_{i\in\ee_j}$ induces the chains $\emptyset\subsetneq \ee_j$ on the $\HH_j$ and $(u_{i})_{i\notin E}$ induces the chain $\emptyset\subsetneq F_1\subsetneq \dots\subsetneq F_\ell$ on $\HH'$.  Thus we have
  \begin{enumerate}
  \item $u_{i_1}=u_{i_2}$ if $i_1,i_2\in\ee_j$ for some $j=1,\dots,r$,
  \item $u_{i_1}\geq u_{i_2}$ if $i_1\in F_{j_1}\setminus F_{j_1-1}$ and $i_2\in F_{j_2}\setminus F_{j_2-1}$ with $j_1\leq j_2$.
  \end{enumerate}
  % As both $\tau$ and $\Star_{\sigma(F_{[m],\bullet})}(v)$ are invariant under translation by the $\bbone_{\ee_j}$, we may assume that $u_i=0$ for all $i\in E$.
  We now show that for $\varepsilon>0$ sufficiently small, $w\coloneqq v+\varepsilon\cdot u$ induces $F_{[m],\bullet}$ on $\GG$.
  For that let $w_{i_1}$ and $w_{i_2}$ be two coordinates of $w$ for $i_1, i_2\in [m]$, $i_1\neq i_2$.  We now consider three cases:

  \noindent
  If $i_1,i_2\in E$, say $i_1\in \ee_{j_1}$ and $i_2\in \ee_{j_2}$ with $j_1\leq j_2$, then we have
  \begin{equation*}
    w_{i_1}=v_{i_1}+ \varepsilon\cdot u_{i_1} = r+1-j_1+ \varepsilon\cdot u_{i_1} \geq r+1-j_2 + \varepsilon\cdot u_{i_2}= v_{i_1}+\varepsilon\cdot u_{i_2}=w_{i_2}.
  \end{equation*}
  If $i_1\in E$, say $i_1\in\ee_j$, and $i_2\notin E$, then we have
  \begin{equation*}
    w_{i_1}=r+1-j + \varepsilon\cdot u_{i_1} > \varepsilon\cdot u_{i_2} = w_{i_2}.
  \end{equation*}
  If $i_1,i_2\notin E$, say $i_1\in F_{j_1} \setminus F_{j_1-1}$ and $i_2\in F_{j_2} \setminus F_{j_2-1}$ for some $j_1\leq j_2$, then we have
  \begin{equation*}
    w_{i_1}=\varepsilon\cdot u_{i_1} \geq \varepsilon\cdot u_{i_2}=w_{i_2}.
  \end{equation*}

  Conversely, let $u\in\Star_{\sigma(F_{[m],\bullet})}(v)$, i.e., $w=v+\varepsilon\cdot u$ induces $F_{[m],\bullet}$ on $\GG$.  Then we have
  \begin{enumerate}
  \item \label{enumitem:relevantCondition1}$w_{i_1}=w_{i_2}$ for all $i_1,i_2\in \ee_j$ and all $j=1,\dots,r$,
  \item $w_{i_1}\geq w_{i_2}$ if either
    \begin{enumerate}
    \item $i_1\in \ee_{j_1}$ and $i_2\in\ee_{j_2}$ with $j_1\leq j_2$,
    \item $i_1\in E$ and $i_2\notin E$,
    \item \label{enumitem:relevantCondition2c} $i_1\in F_{j_1}\setminus F_{j_1-1}$ and $i_2\in F_{j_2}\setminus F_{j_2-1}$ with $j_1\leq j_2$.
    \end{enumerate}
  \end{enumerate}
  By Condition~\eqref{enumitem:relevantCondition1}, the $(u_i)_{i\in\ee_j}$ induce the chains $\emptyset\subsetneq \ee_j$ on $\HH_j$ and, by Condition \eqref{enumitem:relevantCondition2c}, $(u_i)_{i\notin E}$ induces the chain $\emptyset \subsetneq F_1 \subsetneq F_2\subsetneq \dots\subsetneq F_\ell$ on $\HH'$.  Hence $u\in\tau$.

  The ``$\supseteq$'' inclusion follows similarly.  Consider a relative interior point $v\in\mathrm{Relint}(\sigma)$, say $v = \bbone_{\ee_1}+\dots+\bbone_{\ee_1\cup\dots\cup\ee_r}$, and let $\Star_{\sigma(F_{[m],\bullet})}(v)\in \Star_{\Trop(\GG)}(\sigma)$ be a maximal cone that arises from a maximal chain of flats $F_{[m],\bullet}$ extending $F_{E,\bullet}$,
  \begin{equation*}
    F_{[m],\bullet}\colon\;\; \emptyset\subsetneq \ee_1\subsetneq \dots\subsetneq \ee_1\cup\dots\cup\ee_r\subsetneq E\cup F_1\subsetneq E\cup F_2\subsetneq\dots\subsetneq E\cup F_\ell\subsetneq [m].
  \end{equation*}

  By \cref{cor:excisionsAndChains2}, this is also a valid maximal chain of flats on $\HH$ which we can restrict to maximal chains on the connected components of $\HH$: $F_{\HH_j,\bullet}$ on the $\HH_j$ for $j=1,\dots,r$ and $F_{\HH',\bullet}$ on $\HH'$.  This gives rise to a maximal cone $\tau\coloneqq \sigma((F_{\HH_1,\bullet},\dots,F_{\HH_r,\bullet},F_{\HH',\bullet}))\in\Trop(\HH)$.  We now claim that $\tau=\Star_{\sigma(F_{[m],\bullet})}(v)$.  This can be shown using the same steps as before.
\end{proof}

\end{document}
